\documentclass[10pt,a4paper]{amsart}
\usepackage[margin=19mm]{geometry}
\usepackage{amsmath,amssymb,mathtools}
\usepackage[scr=rsfs,cal=euler]{mathalfa}
\usepackage{xfrac}
\usepackage[unicode,hidelinks,bookmarks=false]{hyperref}
\newcommand{\ie}{i.e.\ }
\newcommand{\cf}{cf.\ }
\newcommand{\resp}{respectively\ }
\newcommand{\Subsection}{\subsection}
\providecommand{\nobreakdash}{\nobreak\hbox{-}}
\setlength{\emergencystretch}{2em}
\multlinegap0pt
\usepackage{amssymb}
\usepackage{xcolor}
\newtheorem{lemma}{Lemma}
\newtheorem{corollary}{Corollary}
\newcommand{\AS}{ A$^s$MDS }
\newcommand{\G}{\mathcal{G}}
\title{Projective systems and bounds on the length of codes of non-zero defect}
\author{Tim L. Alderson, Zhipeng Zhang}
\date{Source: arXiv:2504.19325v2 (CC BY 4.0)}
\begin{document}
\maketitle

\noindent\emph{Source and licence.} Projective systems and bounds on the length of codes of non-zero defect. Version arXiv:2504.19325v2 (CC BY 4.0), distributed by arXiv under the Creative Commons Attribution 4.0 International licence, http://creativecommons.org/licenses/by/4.0/, as recorded in the arXiv licence metadata for this exact version and shown on the abstract page https://arxiv.org/abs/2504.19325v2. Full licence notice: \texttt{licenses/arxiv-2504.19325-CC-BY-4.0.txt}.\par\smallskip

\noindent\textit{Editorial scope.} Counted unit: the article's corollary divisors (source0.tex lines 632--639). The article's complete original proof of it is delivered with it (source0.tex lines 640--654, one \texttt{proof} environment of 15 lines). No mathematics is written, repaired or re-derived: the statement and the proof are the article's own lines, in the article's own notation, and no retained line is substituted or re-flowed.  Retained uncounted as the source-local context the proof needs: lines 606--608; lines 617--619.  Nothing was omitted from any retained span, so the package carries no omission record.  No retained line contains a citation, so the wrapper carries no bibliography and no bibliography entry.  No retained line contains a figure environment, an image-inclusion command or drawing code, so no image is delivered and the package declares no asset dependency.  Editorial formatting only, in addition: an AMS \texttt{amsart} preamble that restates the article's own declarations and macros used by the retained lines, namely the environments \texttt{lemma}, \texttt{corollary} on separate counters, together with the standard packages those lines need.  The retained environments print in this document as Lemma 1, Corollary 1 in order of appearance, whereas the article prints the same items with the numbering of its own full document.  The complete native source of the article as distributed in the arXiv e-print for this exact version is bundled unchanged as \texttt{sources/177169/source0.tex}.
	\begin{lemma}\label{lem: full length comb}
		Let $k\ge 3$. If $\G$ is an $[n,k,d]_q$ \AS code with $n=(s+1)(q+1)+k-2$, then for each $0\le j\le k-2$,    $\gamma_j = \frac{{n-j\choose k-1-j}}{{n-d-j \choose k-1-j}}$ is an integer.
	\end{lemma}

	\begin{lemma}\label{lem: alpha beta}
		Let $k\ge 3$. If $\G$ is an $[n,k,d]_q$ \AS code with $n=(s+1)(q+1)+k-3$, then for each $0\le j\le k-3$,  $\alpha_j$, and $\beta_j$ are integers, where $\alpha_j = \frac{{n-j\choose k-2-j}}{{k+s-2-j \choose k-2-j}}$ and  $\beta_j = \alpha_j \cdot  \frac{q(s+1)}{k+s-1-j}$.
	\end{lemma}

	\begin{corollary}\label{cor: prime divisors}
		Let $\G$ be an $[n,k,d]_q$ \AS code.
		\begin{enumerate}
			\item If there exists a prime $p$ with $\prod_{i=1}^{s} (q(s+1)+i) \equiv 0 \mod p$ where $s<p<k+s$, then $n\le (s+1)(q+1)+k-3$. \label{part: 1 cor: prime divisors}
			\item If there exists a prime $p$ with $\prod_{i=1}^{s} (q(s+1)+i) \equiv 0 \mod p$ where $s<p<k+s-1$, then $n\le (s+1)(q+1)+k-4$. \label{part: 2 cor: prime divisors}
			\item If there exists a prime $p$ with $\prod_{i=0}^{s} (q(s+1)+i) \equiv 0 \mod p$, and gcd$(p,q)=1$ where $s+1<p<k+s-1$, then $n\le (s+1)(q+1)+k-4$. \label{part: 3 cor: prime divisors}
		\end{enumerate}
	\end{corollary}

	\begin{proof}
		For part 1, suppose $\G$ is an $[n,k,d]_q$ \AS code with $n=(s+1)(q+1)+k-2$, and let $p$ be a prime with $\prod_{i=1}^{s} (d+i) \equiv 0 \mod p$, $s<p<k+s$. By Lemma \ref{lem: full length comb}, for any $0\le j \le k-1$
		\begin{equation}
			\gamma_j = \frac{{n-j\choose k-1-j}}{{n-d-j \choose k-1-j}} = \frac{(n-j)!\cdot s!}{(d+s)!\cdot (n-j-d)!} =\frac{s!{n-j \choose d}}{(d+1)(d+2)\cdots (d+s)}
		\end{equation}
		is an integer.  Taking $j=n-d-p$, we obtain
		\begin{equation}
			\gamma_j  =\frac{s!{d+p \choose d}}{(d+1)(d+2)\cdots (d+s)} =\frac{s!(d+s+1)(d+s+2)\cdots (d+p)}{p!}.
		\end{equation}
		However, since $p\mid d+a$ for some $1\le a\le s$, $p$ does not divide $d+x$ for $s+1\le x <a+p$, whence $\gamma_j$ is not an integer.
		With reference to Lemma \ref{lem: alpha beta}, the proofs of parts 2 and 3 are entirely similar after observing that
		\[
		\alpha_j = \frac{s!{n-j \choose d+1}}{(d+2)(d+3)\cdots (d+s+1)}, \text{ and } \beta_j = \frac{q\cdot (s+1)!{n-j \choose d}}{(d+1)(d+2)\cdots (d+s+1)}.
		\]
	\end{proof}

\end{document}
